\section{Derived functors on derived categories}
\label{section-derived-functors-classical}

\noindent
In practice derived functors come about most often when given an
additive functor between abelian categories.

\begin{situation}
\label{situation-classical}
Here $F : \mathcal{A} \to \mathcal{B}$ is an additive functor between
abelian categories. This induces exact functors
$$
F : K(\mathcal{A}) \to K(\mathcal{B}), \quad
K^{+}(\mathcal{A}) \to K^{+}(\mathcal{B}), \quad
K^{-}(\mathcal{A}) \to K^{-}(\mathcal{B}).
$$
See Lemma \ref{lemma-additive-exact-homotopy-category}.
We also denote $F$ the composition $K(\mathcal{A}) \to D(\mathcal{B})$,
$K^{+}(\mathcal{A}) \to D^{+}(\mathcal{B})$, and
$K^{-}(\mathcal{A}) \to D^-(\mathcal{B})$ of $F$ with the localization
functor $K(\mathcal{B}) \to D(\mathcal{B})$, etc. This situation leads
to four derived functors we will consider in the following.
\begin{enumerate}
\item The right derived functor of
$F : K(\mathcal{A}) \to D(\mathcal{B})$
relative to the multiplicative system $\text{Qis}(\mathcal{A})$.
\item The right derived functor of
$F : K^{+}(\mathcal{A}) \to D^{+}(\mathcal{B})$
relative to the multiplicative system $\text{Qis}^{+}(\mathcal{A})$.
\item The left derived functor of
$F : K(\mathcal{A}) \to D(\mathcal{B})$
relative to the multiplicative system $\text{Qis}(\mathcal{A})$.
\item The left derived functor of
$F : K^{-}(\mathcal{A}) \to D^{-}(\mathcal{B})$
relative to the multiplicative system $\text{Qis}^-(\mathcal{A})$.
\end{enumerate}
Each of these cases is an example of
Situation \ref{situation-derived-functor}.
\end{situation}

\noindent
Some of the ambiguity that may arise is alleviated by the following.

\begin{lemma}
\label{lemma-irrelevant}
In
Situation \ref{situation-classical}.
\begin{enumerate}
\item Let $X$ be an object of $K^{+}(\mathcal{A})$.
The right derived functor of $K(\mathcal{A}) \to D(\mathcal{B})$
is defined at $X$ if and only if the right derived functor of
$K^{+}(\mathcal{A}) \to D^{+}(\mathcal{B})$ is defined at $X$.
Moreover, the values are canonically isomorphic.
\item Let $X$ be an object of $K^{+}(\mathcal{A})$.
Then $X$ computes the right derived functor of
$K(\mathcal{A}) \to D(\mathcal{B})$
if and only if $X$ computes the right derived functor of
$K^{+}(\mathcal{A}) \to D^{+}(\mathcal{B})$.
\item Let $X$ be an object of $K^{-}(\mathcal{A})$.
The left derived functor of $K(\mathcal{A}) \to D(\mathcal{B})$
is defined at $X$ if and only if the left derived functor of
$K^{-}(\mathcal{A}) \to D^{-}(\mathcal{B})$ is defined at $X$.
Moreover, the values are canonically isomorphic.
\item Let $X$ be an object of $K^{-}(\mathcal{A})$.
Then $X$ computes the left derived functor of
$K(\mathcal{A}) \to D(\mathcal{B})$ if and only if $X$ computes
the left derived functor of $K^{-}(\mathcal{A}) \to D^{-}(\mathcal{B})$.
\end{enumerate}
\end{lemma}

\begin{proof}
Let $X$ be an object of $K^{+}(\mathcal{A})$.
Consider a quasi-isomorphism $s : X \to X'$ in $K(\mathcal{A})$.
By
Lemma \ref{lemma-complex-cohomology-bounded}
there exists quasi-isomorphism $X' \to X''$ with $X''$ bounded below.
Hence we see that $X/\text{Qis}^+(\mathcal{A})$ is cofinal
in $X/\text{Qis}(\mathcal{A})$. Thus it is clear that (1) holds.
Part (2) follows directly from part (1).
Parts (3) and (4) are dual to parts (1) and (2).
\end{proof}

\noindent
Given an object $A$ of an abelian category $\mathcal{A}$ we get a complex
$$
A[0] = ( \ldots \to 0 \to A \to 0 \to \ldots )
$$
where $A$ is placed in degree zero. Hence a functor
$\mathcal{A} \to K(\mathcal{A})$, $A \mapsto A[0]$.
Let us temporarily say that a partial functor is one that is
defined on a subcategory.

\begin{definition}
\label{definition-derived-functor}
In
Situation \ref{situation-classical}.
\begin{enumerate}
\item The {\it right derived functors of $F$} are the partial functors
$RF$ associated to cases (1) and (2) of
Situation \ref{situation-classical}.
\item The {\it left derived functors of $F$} are the partial functors
$LF$ associated to cases (3) and (4) of
Situation \ref{situation-classical}.
\item An object $A$ of $\mathcal{A}$ is said to be
{\it right acyclic for $F$}, or {\it acyclic for $RF$}
if $A[0]$ computes $RF$.
\item An object $A$ of $\mathcal{A}$ is said to be
{\it left acyclic for $F$}, or {\it acyclic for $LF$}
if $A[0]$ computes $LF$.
\end{enumerate}
\end{definition}

\noindent
The following few lemmas give some criteria for the existence of
enough acyclics.

\begin{lemma}
\label{lemma-subcategory-left-resolution}
Let $\mathcal{A}$ be an abelian category. Let
$\mathcal{P} \subset \Ob(\mathcal{A})$ be a subset containing $0$
such that every object of $\mathcal{A}$ is a quotient of an element of
$\mathcal{P}$. Let $a \in \mathbf{Z}$.
\begin{enumerate}
\item Given $K^\bullet$ with $K^n = 0$ for $n > a$
there exists a quasi-isomorphism $P^\bullet \to K^\bullet$
with $P^n \in \mathcal{P}$ and $P^n \to K^n$ surjective
for all $n$ and $P^n = 0$ for $n > a$.
\item Given $K^\bullet$ with $H^n(K^\bullet) = 0$ for $n > a$
there exists a quasi-isomorphism $P^\bullet \to K^\bullet$
with $P^n \in \mathcal{P}$ for all $n$ and $P^n = 0$ for $n > a$.
\end{enumerate}
\end{lemma}

\begin{proof}
Proof of part (1). Consider the following induction hypothesis $IH_n$:
There are $P^j \in \mathcal{P}$, $j \geq n$, with $P^j = 0$ for $j > a$,
maps $d^j : P^j \to P^{j + 1}$ for $j \geq n$, and surjective maps
$\alpha^j : P^j \to K^j$ for $j \geq n$ such that the diagram
$$
\xymatrix{
& &
P^n \ar[d]^\alpha \ar[r] &
P^{n + 1} \ar[d]^\alpha \ar[r] &
P^{n + 2} \ar[d]^\alpha \ar[r] &
\ldots \\
\ldots \ar[r] &
K^{n - 1} \ar[r] &
K^n \ar[r] &
K^{n + 1} \ar[r] &
K^{n + 2} \ar[r] &
\ldots
}
$$
is commutative, such that $d^{j + 1} \circ d^j = 0$ for $j \geq n$,
such that $\alpha$ induces isomorphisms
$\Ker(d^j)/\Im(d^{j - 1}) \to H^j(K^\bullet)$
for $j > n$, and such that $\alpha : \Ker(d^n) \to \Ker(d_K^n)$
is surjective. Then we choose a surjection
$$
P^{n - 1}
\longrightarrow
K^{n - 1} \times_{K^n} \Ker(d^n) =
K^{n - 1} \times_{\Ker(d_K^n)} \Ker(d^n)
$$
with $P^{n - 1}$ in $\mathcal{P}$. This allows us to extend the diagram
above to
$$
\xymatrix{
& P^{n - 1} \ar[d]^\alpha \ar[r] &
P^n \ar[d]^\alpha \ar[r] &
P^{n + 1} \ar[d]^\alpha \ar[r] &
P^{n + 2} \ar[d]^\alpha \ar[r] &
\ldots \\
\ldots \ar[r] &
K^{n - 1} \ar[r] &
K^n \ar[r] &
K^{n + 1} \ar[r] &
K^{n + 2} \ar[r] &
\ldots
}
$$
The reader easily checks that $IH_{n - 1}$ holds with this choice.

\medskip\noindent
We finish the proof of (1) as follows.
First we note that $IH_n$ is true for $n = a + 1$ since
we can just take $P^j = 0$ for $j > a$. Hence we see that
proceeding by descending induction we produce a complex $P^\bullet$
with $P^n = 0$ for $n > a$
consisting of objects from $\mathcal{P}$, and a termwise
surjective quasi-isomorphism $\alpha : P^\bullet \to K^\bullet$ as desired.

\medskip\noindent
Proof of part (2). The assumption implies that the morphism
$\tau_{\leq a}K^\bullet \to K^\bullet$
(Homology, Section \ref{homology-section-truncations})
is a quasi-isomorphism.
Apply part (1) to find $P^\bullet \to \tau_{\leq a}K^\bullet$.
The composition $P^\bullet \to K^\bullet$ is the desired quasi-isomorphism.
\end{proof}

\begin{lemma}
\label{lemma-subcategory-right-resolution}
Let $\mathcal{A}$ be an abelian category. Let
$\mathcal{I} \subset \Ob(\mathcal{A})$ be a subset containing $0$
such that every object of $\mathcal{A}$ is a subobject of an element of
$\mathcal{I}$. Let $a \in \mathbf{Z}$.
\begin{enumerate}
\item Given $K^\bullet$ with $K^n = 0$ for $n < a$
there exists a quasi-isomorphism $K^\bullet \to I^\bullet$
with $K^n \to I^n$ injective and $I^n \in \mathcal{I}$ for all $n$
and $I^n = 0$ for $n < a$,
\item Given $K^\bullet$ with $H^n(K^\bullet) = 0$
for $n < a$ there exists a quasi-isomorphism $K^\bullet \to I^\bullet$
with $I^n \in \mathcal{I}$ and $I^n = 0$ for $n < a$.
\end{enumerate}
\end{lemma}

\begin{proof}
This lemma is dual to Lemma \ref{lemma-subcategory-left-resolution}.
\end{proof}

\begin{lemma}
\label{lemma-subcategory-right-acyclics}
In
Situation \ref{situation-classical}.
Let $\mathcal{I} \subset \Ob(\mathcal{A})$ be a subset with the
following properties:
\begin{enumerate}
\item every object of $\mathcal{A}$ is a subobject of an element of
$\mathcal{I}$,
\item for any short exact sequence $0 \to P \to Q \to R \to 0$ of
$\mathcal{A}$ with $P, Q \in \mathcal{I}$, then $R \in \mathcal{I}$,
and $0 \to F(P) \to F(Q) \to F(R) \to 0$ is exact.
\end{enumerate}
Then every object of $\mathcal{I}$ is acyclic for $RF$.
\end{lemma}

\begin{proof}
Pick $A \in \mathcal{I}$. Let $A[0] \to K^\bullet$
be a quasi-isomorphism with $K^\bullet$
bounded below. Then we can find a quasi-isomorphism
$K^\bullet \to I^\bullet$ with $I^\bullet$ bounded below and
each $I^n \in \mathcal{I}$, see
Lemma \ref{lemma-subcategory-right-resolution}\footnote{By (1)
we see that $\mathcal{I}$ is nonempty. Pick $P$ in $\mathcal{I}$.
Then the short exact sequence $0 \to P \to P \to 0 \to 0$ and assumption (2)
shows that $0$ is in $\mathcal{I}$. Thus the lemma applies.}.
Hence we see that these resolutions are cofinal in the category
$A[0]/\text{Qis}^{+}(\mathcal{A})$. To finish the proof it therefore
suffices to show that for any quasi-isomorphism
$A[0] \to I^\bullet$ with $I^\bullet$ bounded below and $I^n \in \mathcal{I}$
we have $F(A)[0] \to F(I^\bullet)$ is a quasi-isomorphism.
To see this suppose that $I^n = 0$ for $n < n_0$. Of course we may assume
that $n_0 < 0$. Starting with $n = n_0$ we prove inductively that
$\Im(d^{n - 1}) = \Ker(d^n)$ and $\Im(d^{-1})$
are elements of $\mathcal{I}$ using property (2) and the exact sequences
$$
0 \to \Ker(d^n) \to I^n \to \Im(d^n) \to 0.
$$
Moreover, property (2) also guarantees that the complex
$$
0 \to F(I^{n_0}) \to F(I^{n_0 + 1}) \to \ldots \to F(I^{-1}) \to
F(\Im(d^{-1})) \to 0
$$
is exact. The exact sequence
$0 \to \Im(d^{-1}) \to I^0 \to I^0/\Im(d^{-1}) \to 0$
implies that $I^0/\Im(d^{-1})$ is an element of $\mathcal{I}$.
The exact sequence $0 \to A \to I^0/\Im(d^{-1}) \to \Im(d^0) \to 0$
then implies that $\Im(d^0) = \Ker(d^1)$ is an elements of
$\mathcal{I}$ and from then on one continues as before to show that
$\Im(d^{n - 1}) = \Ker(d^n)$ is an element of $\mathcal{I}$
for all $n > 0$. Applying $F$ to each of the short exact sequences
mentioned above and using (2) we observe that $F(A)[0] \to F(I^\bullet)$
is an isomorphism as desired.
\end{proof}

\begin{lemma}
\label{lemma-subcategory-left-acyclics}
In
Situation \ref{situation-classical}.
Let $\mathcal{P} \subset \Ob(\mathcal{A})$ be a subset with the
following properties:
\begin{enumerate}
\item every object of $\mathcal{A}$ is a quotient of an element of
$\mathcal{P}$,
\item for any short exact sequence $0 \to P \to Q \to R \to 0$ of
$\mathcal{A}$ with $Q, R \in \mathcal{P}$, then $P \in \mathcal{P}$,
and $0 \to F(P) \to F(Q) \to F(R) \to 0$ is exact.
\end{enumerate}
Then every object of $\mathcal{P}$ is acyclic for $LF$.
\end{lemma}

\begin{proof}
Dual to the proof of
Lemma \ref{lemma-subcategory-right-acyclics}.
\end{proof}







